\chapter{The Process by Firm Type}\label{ch:iv:the-process-by-firm-type}

A candidate holds an internship offer from a bank that expires in ten days, is between the second and third
rounds at a market maker whose process has taken five weeks so far, and has not heard from a systematic fund
that interviewed her a month ago. Three firm types, three clocks, and a decision due before two of them stop.
Nothing in her preparation for probability questions helps with this one, which is a question about
processes: what each stage is for, how long each firm's stages take, and how to value an offer in hand
against one that may come.

\section{The stages and what each one is for}

Almost every process is the same funnel with different widths. An application is read; a cheap, standardised
test removes most applicants; a screen by telephone or video checks that the test result reflects a person
who can explain it; a final round of several interviews measures the role's core abilities in depth; a
decision is taken; an offer is made. The order is not arbitrary: cheap and valid measurements come first, and
the expensive ones, which cost several senior people's afternoons, come last (One Quant Book~16, chapter~10).

\begin{definition}[Recruiting cycle]\label{def:iv:the-process-by-firm-type:cycle}
A \emph{recruiting cycle}\index{recruiting cycle} is the calendar on which a firm fills an entry-level class:
applications open, assessments, final rounds and offers run in a fixed window, for a start date that is
usually the following year. Hiring outside the cycle, for a specific opening and a start date as soon as the
candidate is free, is lateral hiring.
\end{definition}

\begin{definition}[Hiring committee]\label{def:iv:the-process-by-firm-type:committee}
A \emph{hiring committee}\index{hiring committee} is the group that takes the hiring decision from the written
scores and notes of an \omterm{def:iv:what-each-interview-tests-by-role:loop}{interview loop}, without having interviewed the candidate itself (or with the
interviewers present but not deciding alone).
\end{definition}

A committee exists to make the decision depend on the rubric rather than on the most persuasive interviewer.
Its consequence for the candidate is that everything must survive being written down: a good answer that the
interviewer could not summarise in two lines is a weaker answer than it felt.

\Cref{fig:iv:the-process-by-firm-type:stages} shows the stages by firm type. The differences are in which
extra stages appear (a trading game, a take-home, a track-record review) and in the clock.

\begin{omfigure}
\begin{tikzpicture}[font=\scriptsize,
  st/.style={draw=omDef,rounded corners=2pt,fill=omDef!6,minimum height=0.8cm,text width=1.35cm,inner sep=2pt,align=center},
  ex/.style={draw=omThm,rounded corners=2pt,fill=omThm!8,minimum height=0.8cm,text width=1.35cm,inner sep=2pt,align=center},
  lab/.style={anchor=east,align=right,font=\footnotesize}]
\node[lab] at (-0.1,-0.0) {market maker,\\proprietary firm};
\node[st] at (0.80,-0.0) {apply};
\node[st] at (2.45,-0.0) {timed\\test};
\node[st] at (4.10,-0.0) {screen};
\node[st] at (5.75,-0.0) {technical};
\node[ex] at (7.40,-0.0) {trading\\game};
\node[st] at (9.05,-0.0) {final\\day};
\node[st] at (10.70,-0.0) {offer};
\node[lab] at (-0.1,-1.1) {systematic fund\\(researcher)};
\node[st] at (0.80,-1.1) {apply};
\node[st] at (2.45,-1.1) {screen};
\node[ex] at (4.10,-1.1) {take-\\home};
\node[st] at (5.75,-1.1) {technical};
\node[ex] at (7.40,-1.1) {presen-\\tation};
\node[st] at (9.05,-1.1) {final\\day};
\node[st] at (10.70,-1.1) {offer};
\node[lab] at (-0.1,-2.2) {multi-manager fund\\(portfolio manager)};
\node[st] at (0.80,-2.2) {contact};
\node[ex] at (2.45,-2.2) {track\\record};
\node[st] at (4.10,-2.2) {meetings};
\node[ex] at (5.75,-2.2) {due\\diligence};
\node[ex] at (7.40,-2.2) {terms};
\node[st] at (9.05,-2.2) {committee};
\node[st] at (10.70,-2.2) {offer};
\node[lab] at (-0.1,-3.3) {bank\\(graduate cycle)};
\node[st] at (0.80,-3.3) {apply};
\node[st] at (2.45,-3.3) {video};
\node[st] at (4.10,-3.3) {test};
\node[ex] at (5.75,-3.3) {superday};
\node[st] at (7.40,-3.3) {offer};
\node[st] at (9.05,-3.3) {intern-\\ship};
\node[ex] at (10.70,-3.3) {return\\offer};
\end{tikzpicture}

\omcaption{The stages by firm type, left to right in time. Red boxes are the stages specific to the firm type.
Durations differ widely and are not drawn; a bank's return offer comes a year after the first application.
Stages from the chapter's text and the firms' published descriptions in \cref{dat:iv:the-process-by-firm-type:published}.}\label{fig:iv:the-process-by-firm-type:stages}
\end{omfigure}

\section{Market makers and proprietary firms}

A proprietary trading firm hires traders and researchers mostly into graduate classes and trains them itself,
so its process measures ability more than knowledge: a timed test of arithmetic or probability, a screen, one
or two \omterm{def:iv:what-each-interview-tests-by-role:loop}{technical interviews}, and a final day that often includes a trading game
(\cref{ch:iv:the-final-round-and-trading-games}). Developer processes replace the game by coding and design
interviews. Processes of this kind are short once started, a few weeks from test to decision, and firms
publish how quickly they reply (\cref{dat:iv:the-process-by-firm-type:published}). The candidate's lever is
the first stage: an assessment failed is often not retakeable for many months, and one firm states the
period.

\section{Systematic funds and multi-manager funds}

A systematic fund hiring a quantitative researcher (One Quant Book~17, chapter~17) adds stages that measure
research itself: a take-home dataset or a case study (\cref{ch:iv:research-case-studies-and-take-homes}),
and often a presentation of past work in which the candidate is questioned the way a research review would
question a result.

A multi-manager fund hiring a portfolio manager (One Quant Book~17, chapter~22) runs a different process
altogether. The candidate is not asked for probability puzzles but for a track record: returns, risk, capacity
and the evidence that the returns belong to the candidate and can be moved (One Quant Book~16, chapters 3
and~25, on pods and track record portability). The stages are meetings, due diligence of the record and of
the strategy, negotiation of the terms (payout, drawdown limits, capital), and a committee. Its clock is set
less by the fund than by the candidate's notice period and any non-compete or garden leave at the current
employer (\cref{ch:iv:offers-compensation-and-non-competes}).

\begin{example}[What a track-record review checks]\label{ex:iv:the-process-by-firm-type:record}
A candidate presents three years of monthly returns with an annualised Sharpe ratio of 1.5. The fund asks:
whose capital and whose decisions produced them (the candidate's own book, or a team's); what the gross
exposure, turnover and capacity were; how the returns were computed (net of which costs, marked how); what the
returns were in the months of the three largest market moves; and what the strategy would have returned at the
fund's intended size. With 36 monthly observations, an annualised Sharpe ratio near 1.5 has a standard error
of roughly $\sqrt{(12 + 1.5^2/2)/36} \approx 0.60$ (Lo, 2002, for independent monthly returns; One Quant Book~4, chapter~11): the record is evidence of
skill, not proof of it, and the fund knows it.
\end{example}

\section{Banks and asset managers}

Banks hire graduates on a \omterm{def:iv:the-process-by-firm-type:cycle}{recruiting cycle} that begins about a year before the start date: applications in
the autumn, a video interview and an assessment, a final round the bank itself calls a superday, and an offer
for a summer internship that ends, if it goes well, in a return offer for a full-time place (One Quant
Book~17, chapter~28, on internships, graduate programmes and lateral hires). Asset managers run similar
graduate schemes on a smaller scale. Experienced hires at both come mostly through lateral processes, often
through agency recruiters (\cref{ch:iv:applications}).

\begin{dated}{2026-09}{What firms publish about their process}\label{dat:iv:the-process-by-firm-type:published}
Firms publish little detail, and what they publish is about logistics. Goldman Sachs's student careers page
lists a video interview of about 30 minutes (engineering applicants also take an online coding assessment)
and then a ``Superday'', a series of final-round interviews, typically two to five for campus hires,
depending on the division; its 2027 summer analyst programme was open for applications in September 2026.
Jane Street states that candidates should hear back within a week of their interviews, that it does not
interview at weekends, and that it books and pays for candidates' travel. Optiver's European recruitment FAQ
states that it aims to contact all candidates within seven working days and does not offer assessment resets
or retakes within eight months once an attempt has been started. IMC's US page lists three stages
(application, an assessment for some roles, a series of interviews) and a ten-week summer internship. Two
Sigma's page describes remote video interviews and says the day may be shortened if, after a few interviews,
there is no fit.
\end{dated}

\section{Running several processes at once}

The hook's candidate faces a decision with a deadline, and the decision has the structure of a bet.

\begin{method}[Accept, hold or decline]\label{met:iv:the-process-by-firm-type:accept}
\begin{enumerate}
\item Put a value on the offer in hand, $a$, and on the offer that may come, $b$, in the same units (the method of
\cref{ch:iv:offers-compensation-and-non-competes}; rank them if nothing better).
\item Estimate the chance $q$ that the pending process ends in an offer, from its stage and its pass rates,
and the value $f$ of the fallback if it does not (the next cycle, another firm).
\item Declining to wait is worth $qb + (1-q)f$; accept if $a$ exceeds it. The break-even is
$q^\ast = (a - f)/(b - f)$.
\item Before deciding, ask both firms for time: an extension that lets the pending process finish replaces the
bet by a choice with the facts known. Tell each firm the truth about the other's deadline, without names if
you prefer.
\end{enumerate}
\end{method}

\begin{example}[The hook's decision]\label{ex:iv:the-process-by-firm-type:decision}
Value the bank's internship at 100, a market-maker offer at 130 and the fallback (next year's cycle) at 70. If
the market maker's final round converts two times in five, waiting is worth $0.4 \times 130 + 0.6 \times 70 =
94 < 100$: accept the bank, unless an extension can be had. With an extension, the candidate takes the
market maker's offer if it comes and the bank's otherwise, worth $0.4 \times 130 + 0.6 \times 100 = 112$; the
extension is worth 12, which is why it is always worth asking for.
\end{example}

When offers arrive one at a time and each must be answered before the next is known, the problem is the
classical best-choice problem (Ferguson, 1989): declining the first $r-1$ and then accepting the first offer
better than all earlier ones maximises the chance of taking the best, and for many offers the right $r-1$ is
about $n/e$ (\cref{iq:iv:the-process-by-firm-type:7}). Real processes are kinder, because offers can often be
held for a week and compared; the model's use is to show how much a deadline costs when they cannot.

\section{Worked answers}

\begin{example}[Making three offers land in the same fortnight]\label{ex:iv:the-process-by-firm-type:align}
\emph{``Three processes, from your notes of the firms' published timelines and what their recruiters told you: firm A
takes about 7 weeks from application to decision, firm B about 4, firm C about 10. You want all three decisions in weeks
12 and 13 of your calendar. When do you apply to each, and what can go wrong?''}

\emph{Answer.} Work backwards from week 12: apply to C in week 2, to A in week 5, to B in week 8. Give each a week of
slack, because every stage can slip by a week: then C in week 1, A in week 4, B in week 7, and the decisions spread over
weeks 11 to 13. What goes wrong is the fast firm: B can move from screen to offer in a fortnight and set a deadline of
days, so apply to B last and tell B's recruiter at the screen that you are in other processes, which is true and
ordinary. \emph{Check:} the latest start is the longest process's, $12 - 10 = 2$, and nothing has to begin before
week 1.
\end{example}

\begin{example}[What a committee does with one weak score]\label{ex:iv:the-process-by-firm-type:packet}
\emph{``Your loop's five scores on a 1-to-4 scale are 3, 3, 4, 2 and 3, and the committee's rule is a mean of at least 3
and no score below 2. The 2 carries the note `wrote the function, no tests, did not check the empty input'. What happens,
and what would you have done differently?''}

\emph{Answer.} The mean is $15/5 = 3.0$ and the lowest score is 2, so the rule passes, by the smallest margin it allows:
one more point lost anywhere and the loop fails. Committees read the notes when a packet is on the line, and this note
records something the candidate controlled. Two minutes of testing (the empty input, one element, the example) would
likely have made it a 3 and moved the mean to 3.2. The lesson generalises: the cheapest point in most loops is the one
lost for not checking.
\end{example}

\section{Question bank}

\begin{interviewq}[$\star$ \iqroles{trader, researcher, developer} \iqfirm{any}]\label{iq:iv:the-process-by-firm-type:1}
Why do most processes put a timed test before the screen and the screen before the final round? What would
go wrong in the other order?
\end{interviewq}

\begin{interviewq}[$\star$ \iqroles{trader} \iqfirm{market maker}]\label{iq:iv:the-process-by-firm-type:2}
A market maker's process is a timed arithmetic test, a screen, and a final day with a trading game. You are
strong at games and slow at arithmetic. You have four weeks. How do you split your preparation?
\end{interviewq}

\begin{interviewq}[$\star$ \iqroles{bank} \iqfirm{bank}]\label{iq:iv:the-process-by-firm-type:3}
What happens between the start of a summer internship at a bank and a return offer, and what is being scored
during those ten weeks?
\end{interviewq}

\begin{interviewq}[$\star\star$ \iqroles{trader, researcher} \iqfirm{any}]\label{iq:iv:the-process-by-firm-type:4}
You hold an offer you value at 100 that expires in five days. A process you prefer will end in an offer
worth 130 with probability 0.4; if it does not, your fallback is worth 70. Do you accept? At what probability
would you be indifferent? What is a one-week extension worth to you?
\end{interviewq}

\begin{interviewq}[$\star\star$ \iqroles{researcher, trader} \iqfirm{multi-manager fund}]\label{iq:iv:the-process-by-firm-type:5}
A multi-manager fund is considering you as a portfolio manager and asks for your track record. What will it
verify, what do you prepare, and what will it be sceptical about?
\end{interviewq}

\begin{interviewq}[$\star\star$ \iqroles{risk, bank} \iqfirm{bank}]\label{iq:iv:the-process-by-firm-type:6}
A firm needs 30 graduate hires a year from its internship class. It makes return offers to 60\% of interns,
and 80\% of those accept. How many interns must it take? Which of the two rates would you try to change
first, and how?
\end{interviewq}

\begin{interviewq}[$\star\star\star$ \iqroles{trader, researcher} \iqfirm{any}]\label{iq:iv:the-process-by-firm-type:7}
Five processes will end in five offers, one at a time, in a random order of quality; each must be accepted or
declined on the spot, and you want the best. What rule maximises your chance, and what is that chance? What
happens as the number of offers grows?
\end{interviewq}

\begin{interviewq}[$\star\star\star$ \iqroles{trader} \iqfirm{market maker}]\label{iq:iv:the-process-by-firm-type:8}
You have a final round on Monday at the firm you most want, an offer from a second firm expiring on Friday,
and a third firm asking you for its final round next week. Which conversations do you have, in what order,
and what exactly do you say to each firm?
\end{interviewq}

\begin{omsources}
\item T.~S.~Ferguson, ``Who solved the secretary problem?'', \emph{Statistical Science} 4(3), 1989.
\item A.~W.~Lo, ``The statistics of Sharpe ratios'', \emph{Financial Analysts Journal} 58(4), 2002, 36--52.
\item The firms' careers pages cited in \cref{dat:iv:the-process-by-firm-type:published}, accessed September
2026.
\item One Quant Book~16, chapters 3, 10 and~25; One Quant Book~17, chapters 17, 22 and~28.
\end{omsources}
